% Generated from HPHP_Documentazione_Completa_EDITABILE.docx.
% The phase title is in English; the source content is preserved in Italian.
\section{Phase D --- Evaluation}

\emph{\textbf{FASE D --- Documento 01}}

Team card

Vedi Fase A, Documento 01 --- Team Affordance (Matteo Boscherini,
Alessandro Campedelli, Francesco Maria Fuligni).

\emph{\textbf{FASE D --- Documento 02a}}

Inspection card --- Cognitive Walkthrough

\textbf{Nome della versione valutata}

Version 0 --- v0\_wireframe\_iniziale (Fase C, Documenti 03-05:
Blueprint, Wireframes, Before/After)

\textbf{Dati di contesto dell\textquotesingle ispezione}

\begin{itemize}
\item
  Data e ora: {[}da compilare con la data reale della sessione{]}
\item
  Membri del team coinvolti e ruoli: un membro nel ruolo di Narratore
  (guida il walkthrough seguendo il compito), un membro nel ruolo di
  Ascoltatore/verbalizzatore (annota risposte e criticità), un membro
  nel ruolo di \textquotesingle persona\textquotesingle{} (risponde alle
  domande impersonando Pina in base al suo profilo di Fase B)
\end{itemize}

\textbf{Metodo di ispezione}

\begin{itemize}
\item
  Cognitive Walkthrough
\end{itemize}

\textbf{Compito valutato}

Task equivalente a quello di Fase B (Persona card protagonista): Pina
prenota la visita diabetologica di controllo semestrale, usando le
schermate Wireframe 1→5 (Home, scelta prestazione, autenticazione
leggera, data/ora, riepilogo/conferma).

\emph{Per ciascun passo, il walkthrough risponde alle 4 domande
standard: (1) L\textquotesingle utente cercherà di ottenere
l\textquotesingle effetto giusto? (2) L\textquotesingle utente noterà
che l\textquotesingle azione corretta è disponibile? (3)
L\textquotesingle utente assocerà l\textquotesingle azione corretta
all\textquotesingle effetto desiderato? (4) Se l\textquotesingle azione
corretta viene eseguita, l\textquotesingle utente vedrà che sta facendo
progressi?}

Passo 1 --- Dalla Home, toccare \textquotesingle Prenota una
visita\textquotesingle{} (Wireframe schermata 1)

{\def\LTcaptype{none} % do not increment counter
\begin{longtable}[]{@{}
  >{\raggedright\arraybackslash}p{(\linewidth - 4\tabcolsep) * \real{0.2600}}
  >{\raggedright\arraybackslash}p{(\linewidth - 4\tabcolsep) * \real{0.1200}}
  >{\raggedright\arraybackslash}p{(\linewidth - 4\tabcolsep) * \real{0.6200}}@{}}
\toprule\noalign{}
\endhead
\bottomrule\noalign{}
\endlastfoot
\textbf{Domanda} & \textbf{Risposta} & \textbf{Motivazione} \\
(1) Effetto giusto? & Sì & La visita di controllo è
l\textquotesingle obiettivo esplicito di Pina; l\textquotesingle azione
\textquotesingle Prenota una visita\textquotesingle{} è coerente con
questo obiettivo fin dal titolo. \\
(2) Azione disponibile notata? & Sì & Il pulsante è il più grande e
colorato della schermata, unico elemento con quel colore: alta salienza
visiva anche per chi ha lieve difficoltà visiva (persona Pina,
Fitness=2). \\
(3) Associazione corretta? & Sì & L\textquotesingle etichetta
\textquotesingle Prenota una visita\textquotesingle{} è in linguaggio
naturale e coincide quasi letteralmente con l\textquotesingle obiettivo
di Pina, riducendo il carico interpretativo (Domain=2, Technical=1). \\
(4) Percezione del progresso? & Sì & Il tocco porta immediatamente alla
schermata 2, con transizione visibile e nuovo titolo in alto. \\
\end{longtable}
}

\textbf{Esito del passo}

Successo previsto. Nessuna criticità rilevata su questo passo.

Passo 2 --- Selezionare \textquotesingle Diabetologia (di
nuovo...)\textquotesingle{} (Wireframe schermata 2)

{\def\LTcaptype{none} % do not increment counter
\begin{longtable}[]{@{}
  >{\raggedright\arraybackslash}p{(\linewidth - 4\tabcolsep) * \real{0.2600}}
  >{\raggedright\arraybackslash}p{(\linewidth - 4\tabcolsep) * \real{0.1200}}
  >{\raggedright\arraybackslash}p{(\linewidth - 4\tabcolsep) * \real{0.6200}}@{}}
\toprule\noalign{}
\endhead
\bottomrule\noalign{}
\endlastfoot
\textbf{Domanda} & \textbf{Risposta} & \textbf{Motivazione} \\
(1) Effetto giusto? & Sì & L\textquotesingle opzione evidenziata
corrisponde esattamente alla prestazione ricorrente di Pina. \\
(2) Azione disponibile notata? & Sì & L\textquotesingle opzione è la
prima della lista ed è evidenziata con colore e bordo distinti dalle
altre tre. \\
(3) Associazione corretta? & Parziale & Il testo aggiuntivo
\textquotesingle(di nuovo, come l\textquotesingle ultima
volta)\textquotesingle{} è utile, ma per un utente con bassa competenza
tecnica (Technical=1) il concetto di
\textquotesingle ripetere\textquotesingle{} un\textquotesingle azione
digitale potrebbe non essere immediatamente familiare quanto lo è nel
mondo fisico (es. richiedere la stessa cosa allo sportello). \\
(4) Percezione del progresso? & Sì & Il tocco porta alla schermata 3 con
un nuovo titolo chiaro (\textquotesingle Conferma la tua
identità\textquotesingle). \\
\end{longtable}
}

\textbf{Esito del passo}

Successo previsto, con un\textquotesingle incertezza minore sul punto
(3).

\textbf{Grado assegnato}

Buono (non Eccellente, per l\textquotesingle incertezza sul punto 3).

\textbf{Raccomandazione di design}

{[}W-01{]} Sostituire o affiancare l\textquotesingle icona
\textquotesingle R\textquotesingle{} con un\textquotesingle icona più
universalmente riconoscibile (es. una freccia circolare) e valutare, in
una versione ad alta fedeltà, un micro-test linguistico con utenti reali
sul termine più efficace (\textquotesingle ripeti\textquotesingle,
\textquotesingle prenota di nuovo\textquotesingle, \textquotesingle come
l\textquotesingle ultima volta\textquotesingle).

Passo 3 --- Inserire il codice SMS e confermare (Wireframe schermata 3)

{\def\LTcaptype{none} % do not increment counter
\begin{longtable}[]{@{}
  >{\raggedright\arraybackslash}p{(\linewidth - 4\tabcolsep) * \real{0.2600}}
  >{\raggedright\arraybackslash}p{(\linewidth - 4\tabcolsep) * \real{0.1200}}
  >{\raggedright\arraybackslash}p{(\linewidth - 4\tabcolsep) * \real{0.6200}}@{}}
\toprule\noalign{}
\endhead
\bottomrule\noalign{}
\endlastfoot
\textbf{Domanda} & \textbf{Risposta} & \textbf{Motivazione} \\
(1) Effetto giusto? & Sì & Il testo spiega esplicitamente perché è
richiesto il codice, riducendo l\textquotesingle ambiguità
sull\textquotesingle obiettivo del passo. \\
(2) Azione disponibile notata? & Sì & Un solo campo di inserimento,
isolato visivamente, e un solo pulsante di conferma sotto di esso. \\
(3) Associazione corretta? & Parziale & Pina non ha mai usato
l\textquotesingle autenticazione a due fattori: il concetto di
\textquotesingle codice temporaneo via SMS\textquotesingle{} è nuovo per
lei, anche se più semplice di SPID. Il rischio è che non capisca da dove
prendere il codice (deve uscire dall\textquotesingle app e controllare
gli SMS). \\
(4) Percezione del progresso? & Sì & Dopo la conferma si passa alla
schermata 4 (calendario), con contenuto chiaramente diverso. \\
\end{longtable}
}

\textbf{Grado assegnato}

Scarso, per il punto (3): è il passo più a rischio
dell\textquotesingle intero flusso per il profilo di Pina.

\textbf{Raccomandazione di design}

{[}W-02{]} Aggiungere un\textquotesingle istruzione esplicita e
illustrata (\textquotesingle Controlla i messaggi SMS del tuo telefono,
poi torna qui\textquotesingle) e valutare, in una versione successiva,
un meccanismo di autofill del codice quando tecnicamente disponibile sul
dispositivo, per eliminare il passaggio manuale tra app di messaggistica
e app di prenotazione.

Passo 4 --- Scegliere data e ora (Wireframe schermata 4)

{\def\LTcaptype{none} % do not increment counter
\begin{longtable}[]{@{}
  >{\raggedright\arraybackslash}p{(\linewidth - 4\tabcolsep) * \real{0.2600}}
  >{\raggedright\arraybackslash}p{(\linewidth - 4\tabcolsep) * \real{0.1200}}
  >{\raggedright\arraybackslash}p{(\linewidth - 4\tabcolsep) * \real{0.6200}}@{}}
\toprule\noalign{}
\endhead
\bottomrule\noalign{}
\endlastfoot
\textbf{Domanda} & \textbf{Risposta} & \textbf{Motivazione} \\
(1) Effetto giusto? & Sì & Scegliere quando andare alla visita è un
obiettivo intermedio ovvio per l\textquotesingle utente. \\
(2) Azione disponibile notata? & Sì & Gli slot disponibili sono
evidenziati a colore rispetto a quelli non disponibili; lo slot
\textquotesingle consigliato\textquotesingle{} ha ulteriore enfasi. \\
(3) Associazione corretta? & Sì & Toccare una data/ora per sceglierla è
un\textquotesingle azione fisica coerente con
l\textquotesingle esperienza pregressa di Pina (calendari cartacei,
sportelli). \\
(4) Percezione del progresso? & Sì & La selezione aggiorna visivamente
lo slot scelto prima di passare al riepilogo. \\
\end{longtable}
}

\textbf{Grado assegnato}

Eccellente.

Passo 5 --- Confermare la prenotazione (Wireframe schermata 5)

{\def\LTcaptype{none} % do not increment counter
\begin{longtable}[]{@{}
  >{\raggedright\arraybackslash}p{(\linewidth - 4\tabcolsep) * \real{0.2600}}
  >{\raggedright\arraybackslash}p{(\linewidth - 4\tabcolsep) * \real{0.1200}}
  >{\raggedright\arraybackslash}p{(\linewidth - 4\tabcolsep) * \real{0.6200}}@{}}
\toprule\noalign{}
\endhead
\bottomrule\noalign{}
\endlastfoot
\textbf{Domanda} & \textbf{Risposta} & \textbf{Motivazione} \\
(1) Effetto giusto? & Sì & Il riepilogo motiva chiaramente perché si sta
per confermare. \\
(2) Azione disponibile notata? & Sì & Un solo pulsante primario ben
distinto da \textquotesingle Torna indietro\textquotesingle. \\
(3) Associazione corretta? & Sì & \textquotesingle Conferma
prenotazione\textquotesingle{} è inequivocabile rispetto
all\textquotesingle obiettivo. \\
(4) Percezione del progresso? & Parziale & Il wireframe non mostra
esplicitamente una schermata di
\textquotesingle successo\textquotesingle{} dopo la conferma (fuori
scope del set attuale di schermate): l\textquotesingle utente potrebbe
restare incerto se l\textquotesingle azione sia realmente completata. \\
\end{longtable}
}

\textbf{Grado assegnato}

Buono.

\textbf{Raccomandazione di design}

{[}W-03{]} Aggiungere alla versione successiva una schermata di conferma
esplicita post-azione (\textquotesingle Prenotazione confermata
$\checkmark$\textquotesingle, con riepilogo e opzione per aggiungere al calendario
del telefono), oggi non presente nel set di wireframe.

\textbf{Esito complessivo del Cognitive Walkthrough}

Il flusso è nel complesso ben allineato al profilo di Pina, con due
punti di attenzione: il passo 3 (autenticazione leggera via SMS) resta
il più rischioso, poiché introduce comunque un concetto nuovo per
l\textquotesingle utente; e l\textquotesingle assenza di una schermata
di conferma esplicita a fine flusso (punto 4 del passo 5) rischia di
lasciare l\textquotesingle utente in una condizione di incertezza subito
dopo l\textquotesingle azione più importante del task.

\textbf{Nome della versione ridisegnata (se esistente come esito
dell\textquotesingle ispezione)}

\textbf{DA COMPLETARE --- Da produrre come v1\_dopo\_inspection se il
team decide di implementare le raccomandazioni {[}W-01{]}, {[}W-02{]},
{[}W-03{]} prima di procedere con Formative/Summative test. In tal caso,
aggiornare anche la Design summary card di Fase C (Documento 06).}

\emph{\textbf{FASE D --- Documento 02b}}

Inspection card --- Heuristic Analysis

\textbf{Nome della versione valutata}

Version 0 --- v0\_wireframe\_iniziale (Fase C, Documenti 03-05)

\textbf{Dati di contesto dell\textquotesingle ispezione}

\begin{itemize}
\item
  Data e ora: {[}da compilare{]}
\item
  Membri del team coinvolti: valutazione condotta indipendentemente da
  ciascun membro del team e poi discussa in gruppo, secondo prassi
  standard dell\textquotesingle analisi euristica
\end{itemize}

\textbf{Metodo di ispezione}

\begin{itemize}
\item
  Heuristic Analysis --- 10 euristiche di Nielsen, applicate
  all\textquotesingle intero set di Wireframe (schermate 1-7) e al
  Blueprint
\end{itemize}

{\def\LTcaptype{none} % do not increment counter
\begin{longtable}[]{@{}
  >{\raggedright\arraybackslash}p{(\linewidth - 4\tabcolsep) * \real{0.2600}}
  >{\raggedright\arraybackslash}p{(\linewidth - 4\tabcolsep) * \real{0.1200}}
  >{\raggedright\arraybackslash}p{(\linewidth - 4\tabcolsep) * \real{0.6200}}@{}}
\toprule\noalign{}
\endhead
\bottomrule\noalign{}
\endlastfoot
\textbf{Euristica} & \textbf{Grado} & \textbf{Giustificazione} \\
1. Visibilità dello stato del sistema & Buono & Il
\textquotesingle Prossimo appuntamento\textquotesingle{} in home e il
calendario con slot evidenziati mostrano sempre lo stato corrente; manca
però una conferma esplicita di successo a fine flusso (vedi
{[}W-03{]}). \\
2. Corrispondenza sistema-mondo reale & Eccellente & Linguaggio naturale
(\textquotesingle Prenota una visita\textquotesingle, \textquotesingle I
miei appuntamenti\textquotesingle) al posto di nomi di sistema; metafora
del calendario cartaceo per la selezione data/ora. \\
3. Controllo e libertà dell\textquotesingle utente & Buono &
\textquotesingle Torna indietro\textquotesingle{} presente nel
riepilogo; manca però un tasto di annullamento esplicito durante il
passo di autenticazione leggera (schermata 3). \\
4. Coerenza e standard & Eccellente & Barra di navigazione identica su
tutte le schermate; stile dei pulsanti (colore verde = azione primaria)
coerente in tutto il set. \\
5. Prevenzione degli errori & Scarso & Nella schermata 4 (data/ora) non
è ancora chiaro cosa succeda se l\textquotesingle utente prova a toccare
uno slot non disponibile (grigio): il wireframe non specifica un
feedback di errore/impedimento. \\
6. Riconoscimento piuttosto che ricordo & Eccellente & La prestazione
ricorrente è mostrata direttamente in home; l\textquotesingle utente non
deve ricordare nulla per avviare il task più frequente. \\
7. Flessibilità ed efficienza d\textquotesingle uso & Buono & La
scorciatoia \textquotesingle ripeti\textquotesingle{} per Pina è
efficiente; per un utente esperto come Elena (persona caregiver) manca
però una vista compatta/tabellare pensata per la velocità, oggi la
dashboard \textquotesingle Chi assisto\textquotesingle{} è pensata
principalmente in stile mobile semplice. \\
8. Design estetico e minimalista & Eccellente & Ogni schermata contiene
solo gli elementi strettamente necessari al passo corrente, coerente con
il requisito {[}P-02{]} di Fase B (bassa competenza tecnica di Pina). \\
9. Aiutare a riconoscere e recuperare dagli errori & Scarso & Il set di
wireframe attuale non include una schermata di errore (es. codice SMS
errato, slot non più disponibile al momento della conferma):
comportamento non specificato. \\
10. Guida e documentazione & Buono & La voce \textquotesingle Serve
aiuto?\textquotesingle{} in barra di navigazione è sempre presente; il
suo contenuto interno non è stato ancora progettato in questa iterazione
(fuori scope dei 7 wireframe realizzati). \\
\end{longtable}
}

\textbf{Raccomandazioni di design}

\begin{itemize}
\item
  {[}H-01{]} (euristica 5, 9) Aggiungere alla prossima iterazione due
  schermate di errore esplicite: codice SMS errato/scaduto, e slot non
  più disponibile al momento della conferma --- entrambe con messaggio
  in linguaggio semplice e azione di recupero chiara
\item
  {[}H-02{]} (euristica 3) Aggiungere un\textquotesingle azione di
  annullamento esplicita anche nel passo di autenticazione (schermata
  3), non solo nel riepilogo finale
\item
  {[}H-03{]} (euristica 7) Valutare, in una versione successiva o in una
  variante \textquotesingle desktop\textquotesingle{} per il caregiver,
  una vista più densa e tabellare della dashboard \textquotesingle Chi
  assisto\textquotesingle, più adatta a un utente esperto come Elena che
  gestisce più familiari o più appuntamenti contemporaneamente
\end{itemize}

\textbf{Nome della versione ridisegnata (se esistente come esito
dell\textquotesingle ispezione)}

\textbf{DA COMPLETARE --- Come per il Cognitive Walkthrough: da produrre
come v1\_dopo\_inspection se le raccomandazioni
{[}H-01{]}/{[}H-02{]}/{[}H-03{]} vengono implementate prima delle fasi
di test con utenti reali.}

\emph{\textbf{FASE D --- Documento 03}}

Formative test card

\emph{Una card per ogni versione testata. Questa card riguarda la
Version 0 (o la v1\_dopo\_inspection, se prodotta). Il protocollo
sottostante è pronto per l\textquotesingle esecuzione; le sezioni che
richiedono soggetti reali sono segnalate esplicitamente.}

\textbf{Nome della versione}

{[}da compilare --- Version 0 oppure v1\_dopo\_inspection, secondo
quanto deciso dal team dopo l\textquotesingle Inspection card{]}

\textbf{Setup della sessione di test}

Test in presenza, sugli stessi dispositivi/soggetti già coinvolti nello
User Testing preliminare di Fase A quando possibile (per continuità e
confrontabilità), altrimenti su nuovi soggetti reclutati con gli stessi
criteri. Il wireframe puo essere mostrato tramite prototipo cliccabile
(consigliato, es. esportando le slide come immagini navigabili) o, in
alternativa, tramite simulazione \textquotesingle Mago di
Oz\textquotesingle{} guidata da un membro del team che mostra la
schermata successiva a ogni tocco del soggetto.

\textbf{Metodo di test}

\begin{itemize}
\item
  Discount usability testing (3-5 persone)
\end{itemize}

\textbf{Elenco dei task da testare}

\begin{itemize}
\item
  T1 --- Da Home, prenotare la visita diabetologica di controllo
  ricorrente (schermate 1→5)
\item
  T2 --- Individuare come disdire un appuntamento già fissato (schermata
  6)
\item
  T3 --- (solo per soggetti nel ruolo caregiver) Individuare lo stato
  dell\textquotesingle ultimo appuntamento del familiare assistito
  (schermata 7)
\end{itemize}

\textbf{Metodologia di test}

\begin{itemize}
\item
  Thinking aloud per T1 e T3
\item
  Bottom line data per T2
\end{itemize}

\textbf{Metriche di test}

\begin{itemize}
\item
  Success, Efficiency, Efficacy, Errors, Learnability, Satisfaction
  (SUS)
\end{itemize}

\textbf{Reclutamento dei soggetti (persone reali)}

\textbf{DA COMPLETARE --- Da completare con soggetti reali, con le
stesse caratteristiche indicate per lo User Testing preliminare di Fase
A (segmento A: 65-85 anni, bassa/media competenza digitale; segmento B:
caregiver familiare 45-60 anni). Nome fittizio, età, genere, competenza
tecnica, somiglianza al segmento, relazione col team.}

\textbf{Modalità di reclutamento dei soggetti}

\textbf{DA COMPLETARE --- Da completare, dichiarando onestamente come
sono stati trovati i soggetti.}

\textbf{Motivazione della scelta dei soggetti}

\textbf{DA COMPLETARE --- Da completare con onestà.}

\emph{Le sezioni \textquotesingle Test \#1:
Subject\ldots\textquotesingle{} (una per soggetto reale, con dati
qualitativi/quantitativi, punteggio SUS dal modulo Google del corso,
riflessioni utili per il design) e la sezione conclusiva
\textquotesingle SUS average and variance / Outcomes and reflections
(\textasciitilde300 parole) / Urgency curve / Nome della versione
ridisegnata\textquotesingle{} vanno aggiunte non appena i test saranno
stati condotti. Il protocollo sopra è pronto per essere eseguito.}

\emph{\textbf{FASE D --- Documento 04}}

Summative test card

\emph{Una sola card, sulla versione finale del design ("o la va o la
spacca"). Da eseguire solo dopo aver chiuso il ciclo di test formativi e
le relative modifiche.}

\textbf{Nome della versione finale}

\textbf{DA COMPLETARE --- Da compilare con il nome della versione finale
(main folder di Fase C, eventualmente aggiornata in base agli esiti dei
test formativi).}

\textbf{Setup della sessione di test}

Stesso impianto metodologico del test formativo, preferibilmente con
soggetti nuovi rispetto a quelli già esposti al design nei cicli
precedenti, per evitare effetti di apprendimento che falsino il
risultato conclusivo.

\textbf{Metodo di test}

\begin{itemize}
\item
  Discount usability testing (2-5 persone) oppure Full usability testing
  (20-50 persone), a seconda del tempo disponibile al team; indicare la
  scelta effettiva
\end{itemize}

\textbf{Elenco dei task da testare}

\begin{itemize}
\item
  Stesso set di task del test formativo (T1, T2, T3), per garantire la
  confrontabilità dei risultati prima/dopo
\end{itemize}

\textbf{Metodologia di test}

\begin{itemize}
\item
  Thinking aloud per T1/T3, Bottom line data per T2, coerentemente con
  il test formativo
\end{itemize}

\textbf{Metriche di test}

\begin{itemize}
\item
  Success, Efficiency, Efficacy, Errors, Learnability, Satisfaction
  (SUS) --- obbligatoria la Satisfaction come richiesto dalla card
\end{itemize}

\textbf{Reclutamento dei soggetti (persone reali)}

\textbf{DA COMPLETARE --- Da completare con soggetti reali al termine
del ciclo di test formativi, seguendo gli stessi criteri di
rappresentatività dei segmenti A e B.}

\textbf{Motivazione della scelta dei soggetti}

\textbf{DA COMPLETARE --- Da completare con onestà.}

\emph{Le sezioni \textquotesingle Test \#1:
Subject\ldots\textquotesingle{} e la sezione conclusiva finale (SUS
average and variance, altre metriche quantitative, Conclusions) vanno
aggiunte al termine dell\textquotesingle esecuzione dei test. Questa è
la valutazione conclusiva dell\textquotesingle intero progetto HPHP: va
condotta per ultima, a valle di tutte le altre fasi.}

\textbf{Protocollo di test rapido}

Sessione unica --- 15/20 minuti a persona

\emph{Copre insieme: User Testing preliminare (Fase A) + Formative test
(Fase D)}

\textbf{Team Affordance --- HPHP}

Stampate questo documento e portatelo con voi durante il test.

Prima di iniziare

\textbf{Cosa vi serve}

\begin{itemize}
\item
  Questo documento stampato (o su un secondo dispositivo)
\item
  Un dispositivo con l\textquotesingle app ER Salute installata o il
  sito cupweb.it aperto nel browser (idealmente il telefono/PC della
  persona stessa)
\item
  Le 7 schermate del wireframe v1 pronte da mostrare in sequenza (es.
  \path{phase-c-wireframes-v1.pptx} aperto in modalità presentazione,
  oppure stampate)
\item
  Un modo per prendere appunti (questo documento stesso, o un foglio a
  parte)
\item
  Il link al modulo Google SUS fornito dal corso, aperto e pronto
\end{itemize}

\textbf{Chi coinvolgere}

Bastano 2-3 persone. L\textquotesingle ideale è chi assomiglia il più
possibile ai due segmenti di Fase A: una persona 65-85 anni con
bassa/media familiarità col digitale (segmento A), e/o una persona 45-60
anni che gestisce le pratiche di un familiare anziano (segmento B). Va
benissimo se sono conoscenti o parenti: basta dichiararlo onestamente
nella card, come richiesto dal corso stesso.

\textbf{Ruoli nel team (se siete in due)}

\begin{itemize}
\item
  Conduttore: legge le istruzioni, osserva, non aiuta se non richiesto
  esplicitamente
\item
  Verbalizzatore: prende appunti su questo documento durante il test
\end{itemize}

\textbf{DA LEGGERE AD ALTA VOCE:} \emph{"Grazie per il tuo aiuto. Ti
mostrerò due cose legate a un\textquotesingle app per prenotare visite
mediche, e ti chiederò di provare a fare alcune azioni pensando ad alta
voce, cioè dicendomi cosa stai guardando e cosa stai per fare. Non
c\textquotesingle è nulla di giusto o sbagliato: sto testando
l\textquotesingle app, non te. Se ti blocchi va benissimo, è
un\textquotesingle informazione utile. Possiamo interrompere quando
vuoi."}

Parte 1 --- Il sito/app reale (5 minuti)

\textbf{Tempo stimato: circa 5 minuti}

\emph{Copre: User Testing preliminare card, Fase A (Documento 06a).}

\textbf{DA LEGGERE AD ALTA VOCE:} \emph{"Apri il sito cupweb.it (o
l\textquotesingle app ER Salute) come se dovessi prenotare una visita di
controllo per te stesso/a. Dimmi ad alta voce cosa stai guardando e cosa
faresti."}

\textbf{Task A --- Trovare come prenotare (senza completare la
prenotazione)}

Lascia che la persona esplori per un massimo di 2 minuti. Non
intervenire, annota dove clicca, dove esita, cosa dice.

\textbf{Task B --- Trovare come si disdirebbe un appuntamento}

Chiedi: "Se avessi già una visita prenotata e non potessi più andarci,
dove andresti a guardare per disdirla?" Lascia indicare senza completare
l\textquotesingle azione.

\textbf{Annotazioni --- Task A e B}

{\def\LTcaptype{none} % do not increment counter
\begin{longtable}[]{@{}
  >{\raggedright\arraybackslash}p{(\linewidth - 2\tabcolsep) * \real{0.3500}}
  >{\raggedright\arraybackslash}p{(\linewidth - 2\tabcolsep) * \real{0.6500}}@{}}
\toprule\noalign{}
\endhead
\bottomrule\noalign{}
\endlastfoot
\textbf{Osservazione} & \textbf{Note} \\
Riesce a capire dove cliccare per prenotare? & \\
Momenti di esitazione o confusione (dove, cosa ha detto) & \\
Riesce a individuare la disdetta? & \\
Commenti spontanei rilevanti (citazioni dirette) & \\
\end{longtable}
}

Parte 2 --- Il nuovo design, wireframe v1 (7 minuti)

\textbf{Tempo stimato: circa 7 minuti}

\emph{Copre: Formative test card, Fase D (Documento 03). Mostrate le
schermate del wireframe v1 una alla volta, come se fosse
l\textquotesingle app vera.}

\textbf{DA LEGGERE AD ALTA VOCE:} \emph{"Ora ti mostro una proposta di
come potrebbe diventare questa app. Immagina che sia già installata sul
tuo telefono. Ti chiedo di nuovo di pensare ad alta voce."}

\textbf{Task C --- Prenotare la visita ricorrente (schermate 1→5b)}

Mostra la schermata 1 (Home). Chiedi: "Come prenoteresti la tua prossima
visita di controllo?" Segui il suo percorso mostrando le schermate
corrispondenti a ogni tocco/scelta indicata.

\textbf{Task D --- Disdire un appuntamento (schermata 6)}

Mostra la schermata 6. Chiedi: "E se dovessi disdire, cosa faresti qui?"
--- questo task va lasciato in bottom line data: osserva senza
commentare durante l\textquotesingle esecuzione.

\textbf{Annotazioni --- Task C e D}

{\def\LTcaptype{none} % do not increment counter
\begin{longtable}[]{@{}
  >{\raggedright\arraybackslash}p{(\linewidth - 2\tabcolsep) * \real{0.3500}}
  >{\raggedright\arraybackslash}p{(\linewidth - 2\tabcolsep) * \real{0.6500}}@{}}
\toprule\noalign{}
\endhead
\bottomrule\noalign{}
\endlastfoot
\textbf{Osservazione} & \textbf{Note} \\
Trova subito il pulsante \textquotesingle Prenota una
visita\textquotesingle? & \\
Capisce il passo di autenticazione via SMS (schermata 3)? & \\
Reazione alla schermata di conferma finale (5b) & \\
Riesce a trovare \textquotesingle Disdici\textquotesingle{} senza aiuto?
& \\
Confronto spontaneo con la Parte 1 (dice se gli sembra più
semplice/difficile) & \\
\end{longtable}
}

Parte 3 --- Questionario e chiusura (3-5 minuti)

\textbf{Tempo stimato: circa 3--5 minuti}

\textbf{DA LEGGERE AD ALTA VOCE:} \emph{"Ultima cosa: ti faccio 10
semplici domande a cui rispondere con un punteggio da 1 a 5, riferendoti
alla nuova versione che ti ho appena mostrato. Non c\textquotesingle è
una risposta giusta."}

\textbf{Somministrare il modulo Google SUS del corso (schermate v1). Non
modificare l\textquotesingle ordine né le parole delle prime 10
domande.}

\emph{Promemoria delle 10 domande standard SUS (in italiano), da usare
solo come riferimento --- la somministrazione ufficiale avviene tramite
il modulo Google del corso:}

{\def\LTcaptype{none} % do not increment counter
\begin{longtable}[]{@{}
  >{\raggedright\arraybackslash}p{(\linewidth - 2\tabcolsep) * \real{0.3500}}
  >{\raggedright\arraybackslash}p{(\linewidth - 2\tabcolsep) * \real{0.6500}}@{}}
\toprule\noalign{}
\endhead
\bottomrule\noalign{}
\endlastfoot
\textbf{\#} & \textbf{Domanda} \\
1 & Penso che mi piacerebbe usare questo sistema frequentemente \\
2 & Ho trovato il sistema inutilmente complesso \\
3 & Ho trovato il sistema facile da usare \\
4 & Penso che avrei bisogno del supporto di una persona tecnica per
usare questo sistema \\
5 & Ho trovato che le varie funzioni di questo sistema fossero ben
integrate \\
6 & Ho trovato troppa incoerenza in questo sistema \\
7 & Immagino che la maggior parte delle persone imparerebbe a usare
questo sistema molto rapidamente \\
8 & Ho trovato il sistema molto scomodo da usare \\
9 & Mi sono sentito/a molto sicuro/a nell\textquotesingle usare questo
sistema \\
10 & Ho avuto bisogno di imparare molte cose prima di riuscire ad usare
questo sistema \\
\end{longtable}
}

\textbf{Domande di chiusura (facoltative, in linguaggio libero)}

\begin{itemize}
\item
  C\textquotesingle è qualcosa che ti ha confuso più di altro?
\item
  Cosa cambieresti, se potessi cambiare una sola cosa?
\item
  Ti fideresti a usare questa versione da solo/a, senza chiedere aiuto a
  nessuno?
\end{itemize}

\textbf{DA LEGGERE AD ALTA VOCE:} \emph{"Grazie mille per il tempo che
mi hai dedicato, sei stato/a davvero utile."}

Scheda soggetto --- da compilare dopo ogni sessione

\emph{Compilate una copia di questa pagina per ciascuna persona testata.
I dati confluiranno nelle card ufficiali (Fase A, Documento 06a; Fase D,
Documento 03).}

{\def\LTcaptype{none} % do not increment counter
\begin{longtable}[]{@{}
  >{\raggedright\arraybackslash}p{(\linewidth - 2\tabcolsep) * \real{0.3500}}
  >{\raggedright\arraybackslash}p{(\linewidth - 2\tabcolsep) * \real{0.6500}}@{}}
\toprule\noalign{}
\endhead
\bottomrule\noalign{}
\endlastfoot
\textbf{Campo} & \textbf{Valore} \\
Nome (fittizio, per anonimizzare) & \\
Età & \\
Genere & \\
Background tecnico (1 riga) & \\
Somiglianza al segmento (A: paziente / B: caregiver) & \\
Relazione con il team & \\
Data e ora del test & \\
Assistente del test (membro del team) & \\
Durata effettiva & \\
Task A completato? (sì/no/parziale) & \\
Task B completato? (sì/no/parziale) & \\
Task C completato? (sì/no/parziale) & \\
Task D completato? (sì/no/parziale) & \\
Punteggio SUS totale (dal modulo Google) & \\
Citazione/osservazione più utile per il design & \\
\end{longtable}
}
